\documentclass[11pt,a4paper]{article}

%% PACHETE MATEMATICE
\usepackage{amsmath,amssymb,amsfonts}
\usepackage{amsthm}
\usepackage{mathpazo}
\usepackage{eulervm}
\usepackage{enumerate}
\usepackage{color}
\usepackage{graphicx}
\usepackage[all]{xy}
\usepackage{xcolor}
\usepackage{framed}
\usepackage[unicode]{hyperref}
\setcounter{MaxMatrixCols}{20}
\usepackage{multicol}
%\usepackage[romanian]{babel}


%% ASPECT
\linespread{1.05}
\usepackage[T1]{fontenc}
\usepackage[utf8x]{inputenc}
\usepackage[margin=2cm]{geometry}
% \usepackage[color]{showkeys}
\usepackage{anyfontsize}
\usepackage{titlesec}
\titleformat*{\section}{\large\bfseries}

\usepackage{fancyhdr}
\pagestyle{fancy}
\rhead{\small \color{gray} Notițe de Adrian Manea}
\lhead{\small \color{gray} BCD-Aero, 2017---2018, L. Costache \& M. Olteanu}


\usepackage[autostyle = false, style = german]{csquotes}
\usepackage[romanian]{babel}
\newcommand\qq{\enquote}

%% COMENZILE MELE
\renewcommand*{\proofname}{Dem.:}
\newcommand\bb{\mathbb}
\newcommand\dr{\mathrm}
\newcommand\kal{\mathcal}
\newcommand\fk{\mathfrak}
\newcommand\op{\oplus}
\newcommand\ot{\otimes}
\newcommand\ds{\displaystyle}
\newcommand\id{\indent}
\newcommand\nid{\noindent}
\newcommand\seq{\subseteq}
\newcommand\sneq{\subsetneq}
\newcommand\speq{\supseteq}
\newcommand\spneq{\supsetneq}
\newcommand\rar{\rightarrow}
\newcommand\Rar{\Rightarrow}
\newcommand\xrar{\xrightarrow}
\newcommand\lar{\leftarrow}
\newcommand\Lar{\Leftarrow}
\newcommand\xlar{\xleftarrow}
\newcommand\lrar{\leftrightarrow}
\newcommand\Lrar{\Leftrightarrow}
\newcommand\ideal{\trianglelefteq}
\newcommand\ai{\text{ a.\^{i}. }}
\newcommand\sk{(\textit{Schi\c{t}\u{a}}) }
\newcommand\rhup{\rightharpoonup}
\newcommand\rhdn{\rightharpoondown}
\newcommand\lhup{\leftharpoonup}
\newcommand\lhdn{\leftharpoondown}
\newcommand\vid{\emptyset}
\newcommand\wh{\widehat}
\newcommand\ol{\overline}
\newcommand\NN{\mathbb{N}}
\newcommand\RR{\mathbb{R}}
\newcommand\ZZ{\mathbb{Z}}
\newcommand\QQ{\mathbb{Q}}
\newcommand\CC{\mathbb{C}}

%% STILURILE MELE
\newtheoremstyle{dotless}{}{}{\itshape}{}{\bfseries}{:}{ }{}

\theoremstyle{dotless}
\newtheorem{theorem}{Teorem\u{a}}[section]
\newtheorem{proposition}{Propozi\c{t}ie}[section]
\newtheorem{lemma}{Lem\u{a}}[section]
\newtheorem{corollary}{Corolar}[section]
\newtheorem{exercise}{Exerci\c{t}iu}

\newtheoremstyle{dotlessdr}{}{}{}{}{\bfseries}{:}{ }{}
\theoremstyle{dotlessdr}
\newtheorem{example}{Exemplu}[section]
\newtheorem{definition}{Defini\c{t}ie}[section]
\newtheorem{remark}{Observa\c{t}ie}[section]




\begin{document}

\begin{center}
  {\large\textbf{Seminar 9 \\ Polinom Taylor. Serii Taylor}}
\end{center}

\vspace{2cm}


\section{Formula lui Taylor}
\label{sec:taylor}


Orice funcție cu anumite proprietăți poate fi aproximată cu un polinom.
Rezultatul formal este următorul.

\begin{definition}\label{def:taylor-pol}
  Fie $I \seq \RR$ un interval deschis și $f : I \to \RR$ o funcție de clasă $C^m(I)$.
  Pentru orice $a \in I$, definim \emph{polinomul Taylor} de gradul $n \leq m$ asociat
  funcției $f$ în punctul $a$ prin:
  \[
    T_{n, f, a} (x) = \sum_{k = 0}^n \frac{f^{(k)}(a)}{k!} (x - a)^k.
  \]

  \emph{Restul} (eroarea de aproximare) este definit prin:
  \[
    R_{n, f, a} = f(x) - T_{n, f, a}(x).
  \]
\end{definition}

Primele polinoame (de gradul întîi și al doilea) se numesc, respectiv, \emph{aproximarea %
  liniară} și \emph{pătratică} a lui $f$, în jurul lui $a$.


Acest polinom se regăsește și în \emph{formula lui Taylor}, care ne arată legătura lui
strînsă cu orice funcție.

\begin{theorem}[Formula lui Taylor cu restul Lagrange]\label{thm:taylor-lagr}
  Fie $f : I \to \RR$ o funcție de clasă $C^{n+1}(I)$ și $a \in I$. Atunci, pentru orice
  $x \in I$, există $\xi \in (a, x)$ sau $(x, a)$ astfel încît:
  \[
    f(x) = T_{n, f, a}(x) + \frac{(x - a)^{n+1}}{(n + 1)!} f^{(n + 1)} (\xi).
  \]
\end{theorem}

Așadar, din această teoremă știm mai precis eroarea aproximării unei funcții cu polinomul
Taylor asociat.

Următoarele sînt consecințe ale teoremei:
\begin{enumerate}[(a)]
\item Restul poate fi scris sub \emph{forma Peano}: există o funcție $\omega : I \to \RR$,
  cu $\ds\lim_{x \to a} \omega(x) = \omega(a) = 0$ și restul se scrie:
  \[
    R_{n, f, a}(x) = \frac{(x - a)^n}{n!} \omega(x).
  \]
\item Restul poate fi scris și sub \emph{formă integrală}:
  \[
    R_{n, f, a}(x) = \frac{1}{n!} \int_a^x f^{(n + 1)}(t) (x - t)^n dt;
  \]
\item $\ds\lim_{x \to a} \frac{R_{n, f, a}(x)}{(x - a)^n} = 0$.
\end{enumerate}

Aceste noțiuni pot fi mai departe utilizate pentru a studia \emph{seria Taylor}
asociată unei funcții.

\begin{definition}\label{def:seria-taylor}
  Fie $I \seq \RR$ un interval deschis și fie $f : I \to \RR$ o funcție de clasă
  $C^\infty(I)$. Pentru orice $x_0 \in I$, se definește \emph{seria Taylor} asociată
  funcției $f$ în punctul $x_0$ seria de puteri:
  \[
    T = \sum_{n \geq 0} \frac{f^{(n)}(x_0)}{n!} (x - x_0)^n.
  \]

  Dacă $x_0 = 0$, seria se mai numește \emph{Maclaurin}.
\end{definition}

Importanța seriilor Taylor este dată de rezultatul următor:

\begin{theorem}\label{thm:serie-taylor}
  Fie $a < b$ și fie $f \in C^\infty([a, b])$ astfel încît să existe $M > 0$ cu
  proprietatea că $\forall n \in \NN, \forall x \in [a, b], |f^{(n)}(x) | \leq M$.

  Atunci pentru orice $x_0 \in (a, b)$, seria Taylor a lui $f$ în jurul lui $x_0$
  este uniform convergentă pe $[a, b]$ și suma ei este funcția $f$. Adică avem:
  \[
    f(x) = \sum_{n \geq 0} \frac{f^{(n)}(x_0)}{n!} (x - x_0)^n, \forall x \in [a, b].
  \]
\end{theorem}

%%%%%%%%%%%%%%%%%%%%%%%%%%%%%%%%%%%%%%%%%%%%%%%%%%%%%%%%%%%%%%%%%%%%%%

\section{Exerciții}
\label{sec:exc}

1. Să se dezvolte în serie Maclaurin funcțiile, indicînd și domeniul de convergență:
\begin{enumerate}[(a)]
\item $f(x) = (1 + x)^\alpha, \alpha \in \RR$;
\item $f(x) = \frac{1}{1 + x}$;
\item $f(x) = \sqrt{1 + x}$;
\item $f(x) = \ds\int_0^x \frac{\sin t}{t} dt$;
\item $f(x) = \sin^2 x$;
\item $f(x) = \frac{3}{(1 - x)(1 + 2x)}$;
\item $f(x) = \arcsin x$.
\end{enumerate}

\vspace{1cm}

2. Să se dezvolte în jurul punctului $a \in \RR$ funcțiile:
\begin{enumerate}[(a)]
\item $f: \RR^\ast \to \RR, f(x) = \frac{1}{x}, a = 1$;
\item $f: \RR - \{-2, -1\} \to \RR, f(x) = \frac{1}{x^2 + 3x + 2}, a = -4$;
\item $f: \RR \to \RR, f(x) = \cos^2 x, a = \frac{\pi}{4}$.
\end{enumerate}

\vspace{1cm}

3. Să se determine aproximările liniare și pătratice ale următoarelor
funcții în jurul punctelor indicate:
\begin{enumerate}[(a)]
\item $f(x) = x \ln x$ în jurul punctului $a = 1$;
\item $f(x) = \sqrt[3]{x + 1} \sin x$, în jurul punctului $a = 0$.
\end{enumerate}

\vspace{1cm}

4. Să se arate că seria numerică $\ds\sum_{n \geq 0} \frac{(-1)^n}{3n + 1}$ este
convergentă și să se afle suma ei cu ajutorul seriilor de puteri.

\emph{Indicație:} Considerăm seria $\ds\sum_{n \geq 0} (-1)^n \frac{x^{3n+1}}{3n + 1}$
și luăm $x = 1$. Derivăm termen cu termen și obținem seria pentru $\frac{1}{1 + x^3}$.

\vspace{1cm}

5. Să se afle suma seriilor:
\begin{enumerate}[(a)]
\item $\ds\sum_{n \geq 0} (-1)^n \frac{x^{2n+1}}{2n + 1}$;
\item $\ds\sum_{n \geq 1} \frac{(-1)^{n+1}}{n}$;
\item $\ds\sum_{n \geq 0} \frac{(n + 1)^2}{n!}$;
\item $\ds\sum_{n \geq 1} \frac{n^2(3^n - 2^n)}{6^n}$;
\item $\ds\sum_{n \geq 0} \frac{(-1)^n}{2n + 1}$;
\item $\sum_{n \geq 0} \frac{(-1)^n}{3n + 1} x^{3n+1}$.
\end{enumerate}

\emph{Indicații:}
\begin{enumerate}[(a)]
\item Se folosește seria pentru $\arctan x$, derivată termen cu termen;
\item Folosim seria pentru $e^x + xe^x$, din care o obținem pe cea pentru
  $(x + x^2)e^x$, pe care o derivăm termen cu termen. Pentru $x = 1$, avem
  seria cerută.
\item Se calculează seria $\ds\sum_{n \geq 1} n^2 x^n$, obținută din seria
  pentru $x^n$;
\item Se folosește seria pentru $\arctan x$.
\item
  \[
    \sum_{n \geq 0} \frac{(-1)^n}{3n + 1} = %
    \lim_{x \to 1} \sum_{n \geq 0} \frac{(-1)^n}{3n + 1} x^{3n + 1} = %
    \lim_{x \to 1} \int_0^x \sum_{n \geq 0} (-1)^n x^{3n} dx.
  \]
\end{enumerate}

\vspace{1cm}

6. Să se calculeze $\sqrt[4]{10004}$ cu 4 zecimale exacte.

\emph{Indicație:} Se folosește dezvoltarea binomială:
\[
  \sqrt[4]{10004} = 10 \sqrt[4]{1 + \frac{4}{10000}}.
\]

\vspace{1cm}

7. Să se calculeze cu o eroare mai mică decît $10^{-3}$ integralele:
\begin{enumerate}[(a)]
\item $\ds\int_0^{\frac{1}{3}} \frac{\arctan x}{x} dx$;
\item $\ds\int_0^1 e^{-x^2} dx$;
\item $\ds\int_0^{\frac{1}{2}} \frac{\ln(1 + x^2)}{x} dx$;
\item $\ds\int_0^{\frac{1}{2}} \frac{1 - \cos x^2}{x^2} dx$.
\end{enumerate}

\end{document}
